% Generated draft: requires mathematical and source review.
\documentclass[11pt]{article}
\usepackage[margin=25mm]{geometry}
\usepackage{fontspec}
\setmainfont{Latin Modern Roman}
\usepackage{amsmath,amssymb,mathtools,unicode-math}
\setmathfont{Latin Modern Math}
\usepackage{enumitem,needspace,xurl,hyperref}
\setlength{\parindent}{0pt}
\setlength{\parskip}{5pt}
\setlength{\emergencystretch}{3em}
\newcommand{\N}{\mathbb{N}}
\newcommand{\Z}{\mathbb{Z}}
\newcommand{\Q}{\mathbb{Q}}
\newcommand{\R}{\mathbb{R}}
\newcommand{\C}{\mathbb{C}}
\newcommand{\F}{\mathbb{F}}
\newcommand{\A}{\mathbb{A}}
\newcommand{\PP}{\mathbb{P}}
\newcommand{\E}{\mathbb{E}}
\newcommand{\eps}{\varepsilon}
\newcommand{\dd}{\,\mathrm{d}}
\begin{document}
\section*{UnsolvedMath ID: 20002469}
\textbf{Source problem code:} catalogue entry
\subsection{Definitions and mathematical statement}
Gonzalo Contreras's Question 6 asks whether, ``under the same condition on \(\alpha_L\),'' every Aubry set \(\mathcal A_c\) covers the configuration torus \(T^n\). The preceding condition is not included in this source record. Consequently the literal implication has a missing premise. The aid treats an analytic time-periodic Tonelli Lagrangian and an everywhere positive-definite Hessian of its analytic Mather \(\alpha\)-function; that is an explicit conditional reconstruction, not a recovered quotation of the omitted premise.

For \(L:TT^n\times\mathbb T\to\R\), \(\mathbb T=\R/\Z\), assume fiberwise strict convexity, superlinearity, complete Euler--Lagrange flow, and the indicated analyticity. Identify \(H^1(T^n;\R)\) with \(\R^n\), and define
\[
 \alpha_L(c)=-\inf_\mu\int(L-c\cdot v)\,\dd\mu,
\]
where \(\mu\) ranges over invariant probability measures with finite action. With \(L_c=L-c\cdot v+\alpha_L(c)\), integer-time minimal actions define the Peierls barrier \(h_c(x,y)\), and the time-zero projected Aubry set is \(\mathcal A_c^0=\{x:h_c(x,x)=0\}\). The reconstructed target is \(\mathcal A_c^0=T^n\) for every \(c\); a statement at every time slice needs the corresponding time-periodic interpretation.

The aid reports three scoped consequences. If \(\alpha_L\) is analytic and \(D^2\alpha_L(c)\) is positive definite for every \(c\), its superlinearity makes \(\nabla\alpha_L:\R^n\to\R^n\) a global analytic diffeomorphism, and its Legendre dual
\[
 \beta_L(h)=\sup_c(c\cdot h-\alpha_L(c))
\]
is analytic. The aid invokes Massart--Sorrentino for an affirmative autonomous two-torus case under these hypotheses, without a full theorem citation. Also, an intersection of the lifted phase-space Aubry sets for two distinct classes \(c_0,c_1\) forces \(\alpha_L\) to be affine on \([c_0,c_1]\), hence cannot occur under strict convexity. This means a common velocity-and-time phase point, not merely intersecting projections in the base torus. These results do not establish the arbitrary-dimensional time-periodic implication.

The explicit time-dependent family supplied by the aid is
\[
 L(x,v,t)=\tfrac12(v-b(t))^{\mathsf T}A(v-b(t))
 +\partial_t\phi(x,t)+D_x\phi(x,t)\cdot v,
\]
where \(A\) is a constant symmetric positive-definite matrix, \(b\) is real analytic and one-periodic, and \(\phi\) is real analytic and periodic in space and time. Write \(\bar b=\int_0^1b(t)\,\dd t\). The aid's formulas are
\[
 \alpha_L(c)=c\cdot\bar b+\tfrac12c^{\mathsf T}A^{-1}c,\qquad
 \nabla\alpha_L(c)=\bar b+A^{-1}c.
\]
For lifts \(X,Y\) of \(x,y\) and a positive integer \(N\), the normalized action is
\[
 h_c^N(x,y)=\phi(y,0)-\phi(x,0)
 +\frac1{2N}\min_{\ell\in\Z^n}
 (Y-X+\ell-N(\bar b+A^{-1}c))^{\mathsf T}
 A(Y-X+\ell-N(\bar b+A^{-1}c)).
\]
The quadratic remainder lies between zero and \(\lambda_{\max}(A)n/(8N)\), so
\[
 h_c(x,y)=\phi(y,0)-\phi(x,0),\quad
 \mathcal A_c^0=T^n,
 \qquad p=c+D_x\phi(x,t)
\]
are the barrier, full projection, and dual Aubry graph in the physical cotangent variable \(p=\partial_vL\). These graphs foliate each cotangent fiber as \(c\) varies, with calibrated velocity \(v=b(t)+A^{-1}c\). Nonconstant \(b\) gives genuinely time-dependent coefficients, but the example remains a special drift-and-gauge family. The general analytic time-periodic question is reported open in the literature checked by the aid, not independently checked here.
\subsection{Short English statement}
The source omits the preceding hypothesis on Mather's \(\alpha\)-function. Under the aid's positive-Hessian analytic interpretation, it reports several restricted consequences and an exact time-dependent drift-and-gauge family with full Aubry projection in every class, while retaining the general question as open.
\subsection{Sources}
\begin{itemize}
\item Local registry, permanent UnsolvedMath ID \(20002469\), exact public code \texttt{catalogue entry\!}; canonical location \texttt{aim-pdes-notes.json}, \texttt{notes[80]}. The exact public code is \texttt{catalogue entry}.
\item AIM, \emph{New connections between dynamical systems and PDE's}, Question 6 attributed to Gonzalo Contreras; supplied reference \url{https://aimath.org/WWN/dynpde/dynpde.pdf}. Its phrase ``under the same condition'' points outside the supplied excerpt. No PDF was accessed here. No proposer year, original remarks, or literature citation is supplied.
\item Massart--Sorrentino theorem, invoked by the aid for the autonomous two-torus case; no title, year, journal, or theorem number is supplied.
\item Supplied research attempt 1: \texttt{partial\_result}, run status \texttt{open}, rigor label \texttt{proved}, low novelty confidence; registry status \texttt{partially\_solved}, classification \texttt{PARTIAL-PROGRESS}, created and updated 14 August 2026. The aid is machine generated, not peer reviewed, and requires independent verification before citation.
\end{itemize}
\end{document}
